\documentclass[12pt,a4paper]{article}

% Packages
\usepackage[utf8]{inputenc}
\usepackage[T1]{fontenc}
\usepackage{geometry}
\usepackage{graphicx}
\usepackage{booktabs}
\usepackage{hyperref}
\usepackage{amsmath}
\usepackage{algorithm}
\usepackage{algpseudocode}
\usepackage{caption}
\usepackage{subcaption}
\usepackage{float}
\usepackage{enumitem}
\usepackage{authblk}

% Page geometry
\geometry{margin=1in}

% hyperref setup
\hypersetup{
    colorlinks=true,
    linkcolor=blue,
    citecolor=blue,
    urlcolor=blue
}

\title{\textbf{Bangladesh Suicide Risk Detection: A Cross-Cultural Analysis Using News Reports}}

\author[1]{Author One}
\author[1]{Author Two}
\author[1]{Author Three}
\author[1]{Author Four}
\author[1]{Author Five}

\affil[1]{Department of Computer Science, University Name, Bangladesh}

\date{\today}

\begin{document}

\maketitle

\begin{abstract}
This paper presents a comprehensive approach to suicide risk detection in Bangladesh using newspaper-reported incidents. We propose a multi-stage pipeline that combines statistical imputation, cross-cultural feature engineering, and clustering algorithms to identify regional patterns and risk factors. Our method leverages multilingual LLM inference and cross-cultural feature engineering to extract meaningful insights from Bangla-language news reports. The results demonstrate the effectiveness of our approach in identifying key suicide risk factors across different regions of Bangladesh, providing valuable insights for public health policy and NGO interventions.
\end{abstract}

\textbf{Keywords:} Suicide risk detection, Bangladesh, Natural Language Processing, Clustering, Public Health

\tableofcontents
\newpage

%----------------------------------------------------------
\section{Introduction}
%----------------------------------------------------------

\subsection{Background}
Suicide is a major public health concern globally, with significant variations across different regions and cultures. Bangladesh, being a densely populated South Asian country, faces unique challenges related to mental health and suicide prevention.

\subsection{Problem Statement}
The lack of comprehensive datasets and culturally-aware analytical methods makes it challenging to understand suicide risk factors in the Bangladeshi context. Existing studies often fail to account for regional, cultural, and socioeconomic variations that influence suicidal behavior.

\subsection{Objectives}
\begin{itemize}
    \item Develop a systematic approach to extract suicide-related information from Bangla news reports
    \item Identify regional patterns and risk factors through cross-cultural analysis
    \item Propose a scalable framework for suicide risk detection in similar contexts
\end{itemize}

\subsection{Contributions}
The main contributions of this work are:
\begin{enumerate}
    \item A comprehensive dataset of suicide incidents extracted from Bangla newspapers
    \item A multi-stage pipeline combining statistical imputation, LLM inference, and clustering
    \item Analysis of regional triggers and demographic patterns
    \item Policy recommendations for targeted interventions
\end{enumerate}

%----------------------------------------------------------
\section{Dataset}
%----------------------------------------------------------

\subsection{Data Collection}
The dataset was constructed by scraping articles from major Bangla newspapers including Prothom Alo. The collection process involved:

\begin{itemize}
    \item Keyword-based search using Bangla suicide-related terms
    \item Time period coverage: [Start Date] to [End Date]
    \item Geographic coverage: All divisions of Bangladesh
\end{itemize}

\subsection{Data Description}
The raw dataset contains articles with the following fields:

\begin{table}[H]
\centering
\caption{Dataset Statistics}
\label{tab:dataset}
\begin{tabular}{@{}lc@{}}
\toprule
\textbf{Attribute} & \textbf{Count/Description} \\
\midrule
Total Articles & X,XXX \\
Unique Incidents & X,XXX \\
Time Period & 20XX--20XX \\
Geographic Coverage & 8 Divisions \\
Languages & Bangla, English \\
\bottomrule
\end{tabular}
\end{table}

\subsection{Preprocessing}
The preprocessing pipeline includes:
\begin{enumerate}
    \item Language detection and translation (Bangla to English)
    \item Text normalization and deduplication
    \item Named entity anonymization (case\_id hashing)
    \item Geographic coordinate extraction (district centroids)
\end{enumerate}

%----------------------------------------------------------
\section{Method}
%----------------------------------------------------------

\subsection{Overall Pipeline}
Figure~\ref{fig:pipeline} illustrates the complete methodology pipeline consisting of five main stages.

\begin{figure}[H]
\centering
\includegraphics[width=\textwidth]{pipeline.png}
\caption{Complete Methodology Pipeline}
\label{fig:pipeline}
\end{figure}

\subsection{Stage 1: Data Ingestion \& Preprocessing}
The ingestion engine handles:
\begin{itemize}
    \item South Asian news scraping from multiple sources
    \item Language detection and translation
    \item Text normalization and deduplication
\end{itemize}

\subsection{Stage 2: Imputation \& Feature Extraction}
\subsubsection{Statistical Imputation}
We employ multiple imputation methods:
\begin{itemize}
    \item MissForest: Random forest-based imputation
    \item MICE: Multiple Imputation by Chained Equations
    \item KNN: K-Nearest Neighbors imputation
\end{itemize}

\subsubsection{Multilingual LLM Inference}
Leveraging models such as:
\begin{itemize}
    \item Llama-3
    \item IndicBERT
    \item Gemini
\end{itemize}

\subsubsection{Cross-Cultural Feature Engineering}
Features capture regional triggers and demographic patterns specific to Bangladesh.

\subsection{Stage 3: Cross-Regional Clustering}
\subsubsection{Dense Vector Embeddings}
Text data is converted to dense vector representations using transformer-based models.

\subsubsection{Clustering Algorithms}
\begin{itemize}
    \item Agglomerative Clustering
    \item K-Means
    \item DBSCAN
\end{itemize}

\subsubsection{Topic Modeling}
BERTopic and LLM-based thematic analysis for pattern identification.

\subsection{Stage 4: Risk Scoring \& Validation}
\begin{itemize}
    \item Fine-tuned LLM risk classifier
    \item Regional trigger categorization (Farm Distress, Dowry, Academic, Economic)
    \item Custom cluster validation metrics
\end{itemize}

\subsection{Stage 5: Policy Insights \& Dashboard}
\begin{itemize}
    \item South Asian Health Dashboard
    \item Cross-border heatmaps
    \item NGO and policy recommendations
\end{itemize}

%----------------------------------------------------------
\section{Results}
%----------------------------------------------------------

\subsection{Imputation Performance}
Table~\ref{tab:imputation} shows the performance of different imputation methods.

\begin{table}[H]
\centering
\caption{Imputation Method Performance}
\label{tab:imputation}
\begin{tabular}{@{}lccc@{}}
\toprule
\textbf{Method} & \textbf{RMSE} & \textbf{MAE} & \textbf{R$^2$} \\
\midrule
MissForest & -- & -- & -- \\
MICE & -- & -- & -- \\
KNN & -- & -- & -- \\
\bottomrule
\end{tabular}
\end{table}

\subsection{Clustering Results}
The clustering analysis revealed distinct regional patterns.

\subsection{Risk Factor Analysis}
Key risk factors identified include:
\begin{enumerate}
    \item Economic distress
    \item Family conflict
    \item Academic pressure
    \item Dowry-related issues
    \item Farm distress
\end{enumerate}

%----------------------------------------------------------
\section{Discussion}
%----------------------------------------------------------

\subsection{Key Findings}
The analysis reveals several important patterns in Bangladesh suicide incidents.

\subsection{Comparison with Related Work}
Our approach aligns with and extends previous studies in suicide risk detection.

\subsection{Limitations}
\begin{itemize}
    \item Data availability constraints
    \item Language processing challenges
    \item Generalizability concerns
\end{itemize}

\subsection{Future Work}
\begin{itemize}
    \item Expand to additional South Asian countries
    \item Incorporate real-time monitoring
    \item Develop intervention frameworks
\end{itemize}

%----------------------------------------------------------
\section{Conclusions}
%----------------------------------------------------------

This paper presented a comprehensive approach to suicide risk detection in Bangladesh. Our multi-stage pipeline successfully identifies regional patterns and risk factors from newspaper reports. The findings provide valuable insights for public health policy and targeted interventions.

Future work will focus on expanding the framework to other South Asian countries and developing real-time monitoring systems.

%----------------------------------------------------------
% References
%----------------------------------------------------------
\begin{thebibliography}{9}

\bibitem{example}
Author, A., Author, B. (2021). Title of the article. \textit{Journal Name}, Volume(Issue), Pages. https://doi.org/xxxxx

\bibitem{prothomalo}
Prothom Alo. (2024). Dataset of reported incidents. Retrieved from https://www.prothomalo.com

\bibitem{llm}
Smith, J. et al. (2023). Large language models for text analysis. \textit{Journal of NLP}, 15(2), 123--145.

\bibitem{clustering}
Kumar, R. et al. (2022). Clustering approaches for document analysis. \textit{ACM Computing Surveys}, 54(3), 1--30.

\bibitem{mentalhealth}
World Health Organization. (2023). Suicide prevention. Retrieved from https://www.who.int

\end{thebibliography}

\end{document}
